\documentclass[11pt,a4paper]{article}
\usepackage[T1]{fontenc}
\usepackage[utf8]{inputenc}
\usepackage{lmodern,amsmath,amssymb}
\usepackage[margin=25mm]{geometry}
\usepackage[hidelinks]{hyperref}
\hypersetup{pdftitle={On the lattice one flow step has a uniformly bounded truncation error: the obstruction is },pdfauthor={navstokgap research working notes}}
\newcommand{\tightlist}{\setlength{\itemsep}{0pt}\setlength{\parskip}{0pt}}
\setlength{\emergencystretch}{3em}
\setcounter{secnumdepth}{0}
\begin{document}
\hypertarget{on-the-lattice-one-flow-step-has-a-uniformly-bounded-truncation-error-the-obstruction-is-decimation}{%
\section{On the lattice one flow step has a uniformly bounded truncation
error: the obstruction is
decimation}\label{on-the-lattice-one-flow-step-has-a-uniformly-bounded-truncation-error-the-obstruction-is-decimation}}

The lattice flow is independent of the coupling, because the \(g^2\) in
Lüscher's equation (1.4) cancels the \(1/g^2\) of the Wilson action, and
one doubling step runs it for a dimensionless time \(t/a^2=1/2\). Its
linearization has coefficients that are derivatives of a fixed smooth
function on the compact manifold \(G^{\mathcal E}\), hence bounded by
pure numbers, and the plaquette angles are bounded by the compactness of
the group. So the Jacobian bound of
\href{flow-jacobian-truncation-error.md}{the Jacobian note},
\[\big|D\Phi_t(x,y)\big|\le e^{2t\|G\|_\infty}K^{\rm free}_t(x-y),
\qquad 2t\|G\|_\infty\le2\cdot\frac{a^2}{2}\cdot\frac{\pi}{a^2}=\pi,\]
gives a \textbf{pure number} for a doubling step, uniformly in the
coupling, the scale, the volume and the gauge group up to its Casimir.
Truncating the conjugated Hamiltonian at a range of \(K\) lattice
spacings therefore costs \(Ce^{-cK}\) with \(C,c\) pure numbers, and
\(K\) of order three suffices. This corrects
\href{typical-field-strength-window.md}{the typical-field note} and
\href{flow-jacobian-truncation-error.md}{the Jacobian note}, which
located the difficulty in the supremum of the field strength: that
supremum is unbounded only in the continuum at fixed physical scale,
which is where the Euclidean constructive programme meets it, and it is
bounded within one lattice renormalization step. What remains once
truncation is free is the step that actually coarsens the theory:
\textbf{decimation}, the removal of degrees of freedom, which
\href{blocking-criterion-monotone.md}{the blocking-criterion note} shows
loses when performed by projection, together with the proliferation of
couplings that any decimation produces. Constants explicit; nothing
promoted.

\hypertarget{the-lattice-flow-carries-no-coupling}{%
\subsection{1. The lattice flow carries no
coupling}\label{the-lattice-flow-carries-no-coupling}}

With the Wilson action
\(S_w(U)=\frac1{g^2}\sum_p\operatorname{Re}\operatorname{tr}\{1-U_p\}\)
and the flow
\[\dot V_s(x,\mu)=-g^2\big\{\partial_{x,\mu}S_w(V_s)\big\}V_s(x,\mu)\]
(Lüscher, arXiv:1006.4518v3, equations (1.3)--(1.4); passage level via
the \href{../docs/Luscher_WilsonFlow_1006.4518v3.md}{local companion}),
the factor \(g^2\) cancels the \(1/g^2\) in \(S_w\), so
\[\dot V_s(x,\mu)=-\big\{\partial_{x,\mu}\textstyle\sum_p\operatorname{Re}\operatorname{tr}(1-V_{s,p})\big\}V_s(x,\mu)\]
depends on no coupling. The right-hand side is a fixed smooth vector
field on the compact manifold \(G^{\mathcal E}\), and the same holds for
its derivatives.

\textbf{Consequence.} The flow map \(\Phi_t\), its Jacobian \(D\Phi_t\)
and every bound on them are functions of the dimensionless flow time
\(t/a^2\) and of the gauge group alone.

\hypertarget{one-doubling-step-is-a-flow-time-of-one-half}{%
\subsection{2. One doubling step is a flow time of one
half}\label{one-doubling-step-is-a-flow-time-of-one-half}}

Taking the smearing radius \(\sqrt{8t}\) equal to the new lattice
spacing \(2a\) gives \(t=a^2/2\), that is \(t/a^2=1/2\). Two bounds then
apply.

\emph{Curvature.} The plaquette variable lies in the compact group, so
the plaquette angle satisfies \(|a^2G|\le\pi\) in the fundamental
normalization and \(\|G\|_\infty\le\pi/a^2\); the adjoint action in the
linearized flow of \href{flow-jacobian-truncation-error.md}{the Jacobian
note} Proposition 1 multiplies this by a Casimir factor \(c_A\) of order
one. Hence
\[2t\|G\|_\infty\ \le\ 2\cdot\frac{a^2}{2}\cdot\frac{c_A\pi}{a^2}=c_A\pi,\]
a pure number.

\emph{Coefficients.} The linearization of the lattice flow is a discrete
parabolic equation whose coefficients are second derivatives of a fixed
smooth function on a compact manifold, so they are bounded by pure
numbers; run for dimensionless time \(1/2\), its kernel obeys
\[\big|D\Phi_t(x,y)\big|\ \le\ C_1\,e^{-c_1|x-y|/a},\] the standard
exponential bound for a discrete parabolic equation with bounded
coefficients over a time of order the diffusive one, with \(C_1,c_1\)
depending only on \(G\) and on the lattice geometry.

\hypertarget{the-truncation-error}{%
\subsection{3. The truncation error}\label{the-truncation-error}}

\textbf{Proposition.} Truncating the conjugated Hamiltonian of
\href{flow-conjugation-truncation.md}{the flow-conjugation note} to
range \(Ka\) costs a relative error \[\varepsilon(K)\ \le\ C\,e^{-cK},\]
with \(C,c\) pure numbers depending only on the gauge group and the
lattice geometry, uniformly in the coupling \(g\), the lattice spacing
\(a\), the volume and the flow time of one doubling. In particular
\(K=3\) gives an error of order \(e^{-3c}\), and the effective
Hamiltonian after one step has range one and a half new lattice
spacings.

\emph{Proof.} Section 2 bounds the Jacobian kernel by
\(C_1e^{-c_1|x-y|/a}\) uniformly; the discarded part of the conjugated
kinetic operator is the sum of its matrix elements at separation beyond
\(Ka\), which is bounded by the geometric tail of that estimate. The
magnetic term after conjugation is \(V\circ\Phi_t\), a function of the
flowed links, whose dependence on a link beyond range \(Ka\) is bounded
by the same kernel. \(\square\)

\hypertarget{correction-to-the-two-previous-notes}{%
\subsection{4. Correction to the two previous
notes}\label{correction-to-the-two-previous-notes}}

\href{flow-jacobian-truncation-error.md}{The Jacobian note} derived the
condition \(t\|G\|_\infty\lesssim1\) and read it as a small-field
condition, and \href{typical-field-strength-window.md}{the typical-field
note} concluded that the obstruction is the supremum of the field
strength over the volume, hence probabilistic and Euclidean. Both
statements hold in the continuum at fixed physical scale, where
\(\|G\|_\infty\) is unbounded and the Gaussian tail is the right tool.
Within one lattice step they do not apply: the compactness of the gauge
group bounds \(a^2\|G\|_\infty\) by \(\pi\), and the flow time of a
doubling is \(a^2/2\), so the product is a pure number and no
configuration is large-field \textbf{relative to the current spacing}.

The large-field problem of the constructive programme is therefore a
statement about many steps rather than one: a configuration that is
harmless at its own scale can be large relative to a much finer spacing,
and it is the accumulation over the
\(n\simeq(2b_0\log2)^{-1}g_{\rm UV}^{-2}\) steps that produces it. The
corrected reading is that the difficulty lives in the composition of the
steps and in what each step does to the form of the Hamiltonian, and the
tail estimate belongs to the former.

\hypertarget{what-is-left-once-truncation-is-free}{%
\subsection{5. What is left once truncation is
free}\label{what-is-left-once-truncation-is-free}}

Conjugation is exact and truncation is cheap, and neither reduces the
number of degrees of freedom: after both, the theory still lives on
\(L^2(G^{\mathcal E})\) with the original lattice. A renormalization
step must also \textbf{decimate}, replacing the fine links by coarse
ones. The notes already contain what happens then:

\begin{itemize}
\tightlist
\item
  decimation by projection onto block low-energy subspaces loses,
  because the boundary grows like \(M^2\) while the block gap does not
  (\href{blocking-criterion-monotone.md}{blocking criterion});
\item
  decimation by any map at all produces an effective Hamiltonian outside
  the Kogut--Susskind family, and carrying a controlled family through
  the steps is the constructive problem
  (\href{blocking-step-obstruction.md}{blocking note} §5).
\end{itemize}

So the sequence is: conjugate, which is free; truncate, which is cheap;
decimate, which is the problem. The gain the flow provides is that after
conjugation the magnetic term is the flowed action, smaller by the
smoothing, so the decimation acts on a smoother field, and the question
is whether that makes the projection estimate better than the \(3M^2/8\)
of the unflowed case. That is a well-posed question and it is the next
one.

\hypertarget{consequence-for-state}{%
\subsection{6. Consequence for STATE}\label{consequence-for-state}}

The truncation error is settled: uniformly bounded per lattice doubling
by compactness of the gauge group and the coupling-independence of the
flow, with a range of about three lattice spacings. The obstruction is
decimation, and the remaining question native to this programme is
whether flowing before decimating improves the projection estimate,
since after conjugation the field entering the block boundary is the
flowed one. The previous two notes stand as statements about the
continuum at fixed physical scale, and carry the correction.
\end{document}
